% VI/37 reader bridge -- additive, existing proof text preserved.
\paragraph{Equations define schemes through homogeneous quotients.}
Let \(S=k[X_0,\ldots,X_n]\) with the standard grading, and let
\(I\subseteq S\) be a homogeneous ideal.  The associated projective
closed subscheme is
\[
 Y=\operatorname{Proj}(S/I)\hookrightarrow\operatorname{Proj}S
    =\mathbb P_k^n.
\]
This construction contains more information than the set of common
zeros of \(I\).  In the chart where \(X_i\ne0\), the ring of \(Y\)
is the degree-zero part of the localization:
\[
 \Gamma(Y\cap D_+(X_i),\mathcal O_Y)
      =\bigl((S/I)_{X_i}\bigr)_0.
\]
Thus each projective equation is converted into an ordinary affine
polynomial equation by dividing by an appropriate power of \(X_i\).
The local quotient rings then glue automatically along chart overlaps.

\paragraph{Worked example: a smooth projective conic.}
Take \(S=k[X,Y,Z]\) and \(I=(Y^2-XZ)\).
On \(D_+(X)\), put \(u=Y/X\), \(v=Z/X\).
Dividing by \(X^2\) gives \(v=u^2\), hence
\[
 \bigl((S/I)_X\bigr)_0
      \simeq k[u,v]/(v-u^2)\simeq k[u].
\]
On \(D_+(Z)\), put \(w=Y/Z\) and obtain the ring \(k[w]\).
The opens \(D_+(X)\) and \(D_+(Z)\) cover the conic: a projective
point with \(X=Z=0\) would also have \(Y=0\), which is impossible.
On their intersection, the relation \(XZ=Y^2\) gives
\(w=Y/Z=X/Y=1/u\) (and \(u\) is invertible there).
Consequently the two affine lines glue by inversion, just as in
\(\mathbb P_k^1\).  We obtain a scheme isomorphism
\(C\simeq\mathbb P_k^1\), not merely a bijection of points.

\paragraph{Two schemes with the same underlying point.}
In \(\mathbb P_k^1=\operatorname{Proj}k[X,Y]\), compare
\[
 Y_1=\operatorname{Proj}k[X,Y]/(X),\qquad
 Y_2=\operatorname{Proj}k[X,Y]/(X^2).
\]
Both supports are the point \([0:1]\).  On \(D_+(Y)\), with
\(u=X/Y\), their affine rings are respectively
\[
 k[u]/(u)\simeq k,\qquad k[u]/(u^2).
\]
The first is a reduced point and the second is a nonreduced double
point.  The nilpotent class of \(u\) in the second ring is invisible
in the set of points, but visible in the structure sheaf.
The example illustrates why the language of closed \emph{subschemes}
requires specifying the homogeneous ideal, not just its zero set.

\paragraph{Saturation removes irrelevant-vertex artefacts.}
For an ideal \(I\) in the standard graded ring, define its
saturation by the irrelevant ideal \(S_+=(X_0,\ldots,X_n)\):
\[
 I^{\mathrm{sat}}=(I:S_+^\infty)
   =\{f\in S:\text{some }m\ge0\text{ satisfies }S_+^mf\subseteq I\}.
\]
Saturated ideals are the natural homogeneous equations of closed
subschemes of projective space: the ideals \(I\) and
\(I^{\mathrm{sat}}\) give the same sheaf of ideals on all standard
projective charts.  For a direct calculation, let
\(S=k[X,Y]\) and \(I=(X^2,XY)=X(X,Y)\).
Since \((X,Y)X\subseteq I\), we have
\(X\in I^{\mathrm{sat}}\), and one checks
\(I^{\mathrm{sat}}=(X)\).
On \(D_+(Y)\), both \(I\) and \((X)\) give the local equation
\(u=X/Y=0\), so they describe the same reduced projective point.
In contrast, \((X^2)^{\mathrm{sat}}=(X^2)\): its nilpotent double
point survives projective localization.  Saturation removes the
irrelevant-vertex contribution, \emph{not} genuine embedded or
infinitesimal structure on projective space.

\paragraph{Scheme-theoretic intersection keeps multiplicities.}
For homogeneous ideals \(I,J\subseteq S\), their closed
subschemes meet scheme-theoretically in the quotient by \(I+J\):
\[
 V_+(I)\cap_{\mathbb P^n_k}V_+(J)
       \simeq\operatorname{Proj}\bigl(S/(I+J)\bigr).
\]
In \(\mathbb P_k^2\), intersect the parabola closure
\(V_+(YZ-X^2)\) with the line at infinity \(V_+(Z)\).
The resulting quotient is \(k[X,Y,Z]/(YZ-X^2,Z)\), and on
\(D_+(Y)\) its ring is \(k[x]/(x^2)\).
The intersection is therefore a length-two infinitesimal point.
Its underlying support is still just \([0:1:0]\).  This is a
concrete reason to retain quotient rings throughout projective
intersection calculations.
